% HAFF: Holographic Alaya-Field Framework
% Paper B: Coupling, Records, and the Formation of Stable Relations
% October 2026 working revision.

\documentclass[12pt,a4paper]{article}
\usepackage{arxiv}
\usepackage{amsmath,amssymb,amsfonts,amsthm}
\usepackage{hyperref}
\usepackage{booktabs}
\usepackage{array}
\AtBeginDocument{\setlength{\headheight}{14.5pt}}
\renewcommand{\shorttitle}{Coupling, Records, and Stable Relations}

\newtheorem{theorem}{Theorem}[section]
\newtheorem{proposition}[theorem]{Proposition}
\newtheorem{lemma}[theorem]{Lemma}
\newtheorem{corollary}[theorem]{Corollary}
\newtheorem{definition}[theorem]{Definition}
\newtheorem{remark}[theorem]{Remark}
\newtheorem{conjecture}[theorem]{Conjecture}

\newcommand{\Hil}{\mathcal H}
\newcommand{\B}{\mathfrak B}
\newcommand{\E}{\mathbb E}
\newcommand{\Gin}{\Gamma_{\mathrm{in}}}
\newcommand{\Gout}{\Gamma_{\mathrm{out}}}
\newcommand{\ket}[1]{|#1\rangle}
\newcommand{\bra}[1]{\langle #1|}
\newcommand{\proj}[1]{|#1\rangle\!\langle #1|}
\newcommand{\Tr}{\operatorname{Tr}}

\title{Coupling, Records, and Stable Relations:\\
Formation, Protected Content and the Lifetime of a Geometry}

\author{
  Sidong Liu, PhD \\[0.6ex]
  iBioStratix Ltd \\[0.8ex]
  \texttt{sidongliu@hotmail.com}
}

\date{October 2026}

\hypersetup{pdftitle={Coupling, Records, and Stable Relations: Formation, Protected Content and the Lifetime of a Geometry},pdfauthor={Sidong Liu}}
\begin{document}
\emergencystretch=2em
\raggedbottom
\hbadness=5000
\vbadness=5000

\maketitle

\begin{abstract}
A physical coupling determines which distinctions decay, which can be recorded and which remain available for action. For finite systems with commuting Hermitian noise operators, we derive the fixed block algebra and its long-time pinching, then prove a finite-window selection theorem: a separated set of decoherence rates selects at most one block partition, and its algebra contains every stable block algebra in the supplied joint eigenbasis. Finite-time channels still have full operator range. Block-constant recording couplings report only sector labels while protecting the quantum content within each sector; leakage through independent fragments costs coherence multiplicatively, giving a sharp $R^{-1/2}$ symmetry threshold for $R$ records. Geometry read from mutual information receives both sector and protected-content contributions, which we separate by an exact classical--quantum chain identity. Neither contribution is confined to the noise hypergraph: disconnected decoherence-free sectors can preserve an initially shared logical Bell pair. We therefore distinguish preservation from formation. A collective recording interaction creates a flagged protected Bell branch from a product input, with every branch and the nonlocal recording cost retained; a physical Hamiltonian link creates a deterministic Bell pair between two noise-disconnected protected sectors. For the GHZ state, collective and local couplings yield a triangle, one edge or no edge under one reconstruction rule, with an explicit lifetime under symmetry breaking. Beyond commuting noise, complete finite-time record laws and risks vary continuously even when the exact fixed algebra jumps. These results identify coupling-selected storage, record-conditioned formation and finite operational capability without equating them. The ambient carrier, available interactions and reconstruction dictionary remain physical inputs.
\end{abstract}

\section{Introduction}
\label{sec:intro}

A stable distinction is a physical achievement. It must survive the interactions that would erase it, be readable through an available instrument, and last long enough to matter. An algebra of invariant observables describes one part of that achievement; it does not specify a reader, a preparation history or a set of available gates.

Paper~A \cite{Liu2026PaperA} shows why the distinction matters for geometry. Channels generating the same observable algebra can resolve different states and produce non-isometric finite correlation geometries. Its local bit-flip family has a Stinespring realization and a finite fragment-readout certificate. Here we ask a complementary question: what structure a coupling protects, how actual records and interactions form usable relations within that structure, and when those relations lose their support.

Commuting Hermitian noise gives an exact starting point. The fixed observables are the commutant of the coupling operators, familiar from einselection and noiseless codes \cite{Zurek1982,Zurek2003,ZanardiRasetti1997,KnillLaflammeViola2000}. The infinite-time channel pinches onto this algebra. At finite times its operator range is full, so a protected block algebra is additional stability data rather than the range of the actual channel. A window and a tolerance give a unique block partition when intra-block and inter-block rates are separated (Theorem~\ref{thm:window}). The maximality statement concerns block algebras in the specified joint eigenbasis, not every rotated algebra or every useful finite task.

Recording adds a different restriction. A coupling constant on each block can record its label without learning its quantum content. Independent fragments that resolve two states within a block destroy their coherence according to a multiplicative visibility budget (Theorem~\ref{thm:leak}). The centre and the blocks thus perform different jobs within one stable structure. They are not disjoint sources of geometry: classical sector weights can sustain mutual information even when every block is one-dimensional. Section~\ref{sec:sector} gives the exact decomposition needed to distinguish the contributions.

Selection alone also says too little about formation. Two disconnected collective-noise components can preserve a logical Bell pair shared at the start. Conversely, local independent processes cannot create cross-component correlations from independent inputs. Section~\ref{sec:formation} supplies the missing process distinction and two explicit formation mechanisms. One retains a collective measurement flag and all its branches; the other uses an actual cross-component Hamiltonian link. Both generate relations from product inputs, and both expose the interaction resources responsible for them. Foundation~III \cite{PaperFIII} supplies a further mechanism in which acquired records select an encoder and decoder for the same pre-existing quantum content. Storage, formation and learned use then have separate, testable meanings.

The GHZ calculation in Section~\ref{sec:ghz} shows how a coupling changes a finite correlation geometry and how weak symmetry breaking ends its plateau. For noncommuting perturbations, Section~\ref{sec:noncommuting} bounds complete adaptive record laws and task risks without requiring a stable algebra. A transverse qubit exhibits a diverging record lifetime despite a scalar exact fixed algebra.

The supplied Hilbert space, coupling permissions and site factorization are retained throughout. Consciousness~I \cite{PaperCI} selects a Virasoro fixed net inside a given current-net representation; Foundation's hulls concern closure under declared operations. Neither result derives the ambient carrier. Our formation results address how relations and usable protected sectors arise on that carrier. The stronger question of effective subsystem formation returns in Section~\ref{sec:conj}, with a process-sensitive conjecture rather than a noise-only hypergraph rule.

\section{Reductions selected by a coupling}
\label{sec:selected}

\subsection{Setting}

Let $\Hil$ be a finite-dimensional system Hilbert space. The environment couples to the system through Hermitian operators $A_1,\dots,A_K$ with $[A_k,A_l]=0$, at rates $\gamma_k>0$. In the Markovian limit the reduced dynamics is generated by
\begin{equation}
\mathcal L(\rho)=\sum_{k=1}^K\gamma_k\Bigl(A_k\rho A_k-\tfrac12\{A_k^2,\rho\}\Bigr),\qquad \Phi_t=e^{t\mathcal L}.
\label{eq:L}
\end{equation}
A system Hamiltonian that commutes with every $A_k$ only rotates states inside the blocks defined below; we work in the corresponding interaction picture and return to Hamiltonians that do not commute in Remark~\ref{rem:hamiltonian}. Section~\ref{sec:records} treats the same couplings in an exact unitary model with environmental fragments.

Choose an orthonormal basis $\{\ket{x}:x\in\mathcal X\}$ of joint eigenvectors, $A_k\ket x=a_k(x)\ket x$, and put
\begin{equation}
D(x,y)=\tfrac12\sum_k\gamma_k\bigl(a_k(x)-a_k(y)\bigr)^2=\|u(x)-u(y)\|^2,
\qquad u(x)=\bigl(\sqrt{\gamma_k/2}\;a_k(x)\bigr)_{k=1}^K\in\mathbb R^K .
\label{eq:D}
\end{equation}
We call $u(\mathcal X)\subset\mathbb R^K$ the rate image of the basis; decoherence rates are squared Euclidean distances in it. Write $x\approx y$ when $a_k(x)=a_k(y)$ for every $k$, let $\pi_\infty$ be the partition of $\mathcal X$ into $\approx$-classes $a$, and let $P_a$ be the projection onto $\Hil_a=\operatorname{span}\{\ket x:x\in a\}$. For any partition $\pi$ of $\mathcal X$ into classes $C$, put
\[
\B_\pi=\bigoplus_{C\in\pi}\mathcal B(\Hil_C),\qquad
\E_\pi(\rho)=\sum_{C\in\pi}P_C\rho P_C .
\]
$\B_\pi$ is a unital $*$-algebra, $\E_\pi$ is a channel (the pinching), and $\E_\pi$ is its own dual, the conditional expectation onto $\B_\pi$.

\begin{proposition}[The reduction selected by a commuting coupling]
\label{prop:selected}
\begin{enumerate}
\item $\Phi_t$ acts as a Schur multiplier in the eigenbasis: $\Phi_t(\rho)_{xy}=e^{-tD(x,y)}\rho_{xy}$, and $\Phi_t^*=\Phi_t$.
\item The observables fixed by every $\Phi_t^*$ form the algebra $\B_{\pi_\infty}=\{A_1,\dots,A_K\}'$.
\item $\Phi_t\to\E_{\pi_\infty}$ as $t\to\infty$, and the algebra generated by the dual range of $\E_{\pi_\infty}$ is $\B_{\pi_\infty}$. In the terminology of Paper~A, $(\B_{\pi_\infty},\E_{\pi_\infty})$ is a minimally compatible coarse-graining structure.
\item For every finite $t$, $\Phi_t$ is invertible on the operator space, so the algebra generated by its dual range is $\mathcal B(\Hil)$.
\end{enumerate}
\end{proposition}

\begin{proof}
(1) In the eigenbasis, $(A\rho A-\frac12\{A^2,\rho\})_{xy}=(a(x)a(y)-\frac12a(x)^2-\frac12a(y)^2)\rho_{xy}=-\frac12(a(x)-a(y))^2\rho_{xy}$, so $\mathcal L$ multiplies the $(x,y)$ entry by $-D(x,y)$. The multiplier is real and symmetric, hence $\Phi_t^*=\Phi_t$. (2) $X$ is fixed for all $t$ exactly when $X_{xy}=0$ whenever $D(x,y)>0$, that is, whenever $x\not\approx y$. An operator commutes with every diagonal $A_k$ exactly when it has no entries between basis vectors with different eigenvalue tuples. (3) Entries with $D(x,y)>0$ tend to zero and the others are constant. The dual range of $\E_{\pi_\infty}$ is $\B_{\pi_\infty}$, which is already an algebra. (4) All multipliers $e^{-tD(x,y)}$ are nonzero.
\end{proof}

Part (2) is the elementary commuting case of the known structure of fixed points of quantum dynamical semigroups \cite{Frigerio1978,BlumeKohout2010}. For this series the consequence matters more than the statement: the given coupling determines both its channel and its invariant algebra. Part (4) is Paper~A's result that an algebra label does not record resolution, now in dynamical form. Before the long-time limit the channel generates everything, and the structure that a coupling makes accessible has to be identified by how well observables survive over time.

The objects of Paper~E \cite{Liu2026PaperE} can be located in this class. Its environmental fixed set $\operatorname{Fix}(\mathcal E)$ is $\B_{\pi_\infty}$. In Paper~E's collision model, with a qubit pointer, the pointer algebra equals the fixed algebra of the induced dephasing and contains every effect a fragment can report. The centre of $\B_{\pi_\infty}$ is generated by the block projections $P_a$; it is commutative, and Section~\ref{sec:records} shows that it is what the environment records, so it serves as Paper~E's pointer or record algebra. The inclusion of the record algebra in the fixed set is thus a conclusion of the model, as Paper~E requires of any such overlap. The fixed set is strictly larger than its centre as soon as some block has dimension greater than one; a nontrivial two-sector qubit pointer has only one-dimensional blocks. Such blocks are decoherence-free subspaces \cite{ZanardiRasetti1997,KnillLaflammeViola2000}. They retain quantum content; sector weights can retain classical correlations as well.

\begin{remark}[Hamiltonians that do not commute with the couplings]
\label{rem:hamiltonian}
If the system Hamiltonian $H$ has a part $H_{\mathrm{off}}$ connecting different blocks, $\B_{\pi_\infty}$ need not remain exactly invariant. Strong dephasing can then confine the dynamics to the blocks, with leakage on a time scale of order $\Gout/\|H_{\mathrm{off}}\|^2$, where $\Gout$ is the smallest between-block rate; this is the mechanism of quantum Zeno subspaces \cite{FacchiPascazio2002}. Section~\ref{sec:noncommuting} gives an exact two-level lifetime and a finite-time stability bound without a commutativity assumption. Paper~D's chain with a transverse field requires its own many-body estimate.
\end{remark}

\section{Selection over a finite window}
\label{sec:window}

\begin{definition}[Stable algebra on a window]
\label{def:stable}
A unital $*$-subalgebra $\B\subseteq\mathcal B(\Hil)$ is $\varepsilon$-stable on $[0,T]$ if
\[
\sup_{0\le t\le T}\ \sup_{X\in\B,\ \|X\|\le1}\ \|\Phi_t^*(X)-X\|\le\varepsilon .
\]
\end{definition}

The diagnostic tests the whole unit ball of an algebra, rather than selected states or a list of slow observables. For a partition $\pi$ define the width of a class $C$ in the direction $k$, $w_k(C)=\max_{x\in C}a_k(x)-\min_{x\in C}a_k(x)$, and the two rates
\begin{equation}
\Gin(\pi)=\max_{C\in\pi}\ \tfrac12\sum_k\gamma_kw_k(C)^2,
\qquad
\Gout(\pi)=\min\{D(x,y):\ x,y\text{ in different classes of }\pi\}.
\label{eq:rates}
\end{equation}
The minimum over an empty set is $+\infty$, so a single-block partition has $\Gout=+\infty$. If every coupling is scalar, the full algebra is fixed and selected for every window.
In the rate image, $\Gin(\pi)$ is the largest squared diagonal of the axis-parallel boxes enclosing the classes, and $\Gout(\pi)$ is the smallest squared distance between points of different classes.

\begin{theorem}[Selection by a window]
\label{thm:window}
Let $0<\varepsilon<1$ and $T>0$.
\begin{enumerate}
\item For every $X\in\B_\pi$ and $t\ge0$, $\|\Phi_t^*(X)-X\|\le t\,\Gin(\pi)\|X\|$. For every state $\rho\in\B_\pi$, $\|\Phi_t(\rho)-\rho\|_1\le t\,\Gin(\pi)$. In particular $\B_\pi$ is $\varepsilon$-stable on $[0,T]$ when $T\Gin(\pi)\le\varepsilon$.
\item If $1-e^{-T\Gout(\pi)}>\varepsilon$, no unital $*$-algebra that contains $\B_\pi$ properly is $\varepsilon$-stable on $[0,T]$.
\item Call $\pi$ selected by $(T,\varepsilon)$ if
\begin{equation}
T\,\Gin(\pi)\le\varepsilon<1-e^{-T\Gout(\pi)} .
\label{eq:gap}
\end{equation}
At most one partition is selected by $(T,\varepsilon)$. If $\pi$ is selected by $(T,\varepsilon)$ and $\pi'$ by $(T',\varepsilon)$ with $T\le T'$, then $\pi'$ refines $\pi$, so $\B_{\pi'}\subseteq\B_\pi$.
\item For $T>-\ln(1-\varepsilon)/\Gout(\pi_\infty)$ the selected partition exists and is $\pi_\infty$, so the selected algebra is the fixed algebra of Proposition~\ref{prop:selected}.
\end{enumerate}
\end{theorem}

\begin{proof}
(1) $\B_\pi$ is invariant under $\mathcal L$, and on it
$\mathcal L(X)=-\frac12\sum_k\gamma_k[A_k,[A_k,X]]$.
On a class $C$, replacing $A_k$ by $A_k-m_k$ with $m_k$ the midpoint of its eigenvalues on $C$ does not change the commutator and gives an operator of norm $w_k(C)/2$; hence $\|[A_k,X_C]\|\le w_k(C)\|X_C\|$ and $\|[A_k,[A_k,X_C]]\|\le w_k(C)^2\|X_C\|$. Since $\|X\|=\max_C\|X_C\|$, $\|\mathcal L(X)\|\le\Gin(\pi)\|X\|$. Then $\Phi_t^*(X)-X=\int_0^t\Phi_s^*(\mathcal L(X))\,ds$ and $\Phi_s^*$ is a contraction in operator norm. For a state $\rho\in\B_\pi$ and $\|Y\|\le1$, $\Tr[(\Phi_t(\rho)-\rho)Y]=\Tr[\rho(\Phi_t^*-\mathrm{id})(\E_\pi Y)]$, because $\rho=\E_\pi\rho$, $\E_\pi$ is self-dual and commutes with the Schur multiplier $\Phi_t^*$; since $\E_\pi Y\in\B_\pi$ and $\|\E_\pi Y\|\le1$, the first statement bounds this by $t\Gin(\pi)$.

(2) Let $\B\supsetneq\B_\pi$ and $X\in\B\setminus\B_\pi$. Some entry $X_{x'y'}$ with $x'\in C$, $y'\in C'$, $C\ne C'$, is nonzero. For any $x\in C$, $y\in C'$ the matrix units $E_{xx'}$ and $E_{y'y}$ lie in $\B_\pi\subseteq\B$, so $E_{xx'}XE_{y'y}=X_{x'y'}E_{xy}\in\B$. Then $\|\Phi_T^*(E_{xy})-E_{xy}\|=1-e^{-TD(x,y)}\ge1-e^{-T\Gout(\pi)}>\varepsilon$.

(3) Let $\pi$ and $\pi'$ be selected by $(T,\varepsilon)$ and $(T',\varepsilon)$, $T\le T'$. If $x,y$ lie in one class of $\pi'$, then $D(x,y)\le\Gin(\pi')\le\varepsilon/T'\le\varepsilon/T$. If $x,y$ lay in different classes of $\pi$, then $D(x,y)\ge\Gout(\pi)>-\ln(1-\varepsilon)/T\ge\varepsilon/T$. Hence every class of $\pi'$ lies in a class of $\pi$. For $T=T'$ the same argument with the roles exchanged gives $\pi=\pi'$.

(4) $\Gin(\pi_\infty)=0$ and $\Gout(\pi_\infty)>0$, so $\pi_\infty$ satisfies \eqref{eq:gap} for such $T$; uniqueness is part (3).
\end{proof}

A window and a tolerance therefore pick out at most one algebra of the form $\B_\pi$, and they pick one out when the rate image of the basis splits into clusters that are small compared with $\varepsilon/T$ and far apart compared with $-\ln(1-\varepsilon)/T$. This is a separation of time scales, the mechanism behind metastability in open quantum systems \cite{Macieszczak2016}; for commuting couplings it can be stated exactly. The selected algebra is moreover the largest $\varepsilon$-stable algebra of block form: if $\B_\sigma$ is $\varepsilon$-stable on $[0,T]$, every matrix unit $E_{xy}$ with $x,y$ in one class of $\sigma$ satisfies $1-e^{-TD(x,y)}\le\varepsilon$, so $x$ and $y$ lie in one class of the selected partition $\pi$, and $\B_\sigma\subseteq\B_\pi$. As the window lengthens, the selected algebra shrinks in discrete steps located at the gaps of the rate spectrum and ends at the fixed algebra. In the order of Paper~A this is a chain of further coarse-grainings, since $\E_{\pi'}=\E_{\pi'}\circ\E_\pi$ when $\pi'$ refines $\pi$.

Without a gap there is no selected algebra at that window. Selection uses the sufficient bound of part~(1), so between that bound and the necessary bound of part~(2) a unique maximal stable block algebra can still exist; but the window may also fail to determine one, and two incomparable maximal stable block algebras can occur. The set of $\varepsilon$-stable matrix units can then be an operator system whose generated algebra contains unstable elements: if $D(x,y)$ and $D(y,z)$ are small, $D(x,z)$ can still be as large as $(\sqrt{D(x,y)}+\sqrt{D(y,z)})^2$, and a chain of small steps can connect distant points of the rate image. Theorem~\ref{thm:window} is a sufficient selection certificate. Failure of its gap condition does not by itself decide existence or uniqueness of a maximal stable algebra. Nor does stability of a block algebra supply the instruments needed to use every element as a gate.

\subsection{Continuous capabilities when the fixed algebra jumps}
\label{sec:noncommuting}

A weak coupling can change the exact commutant discontinuously. It cannot
change every finite experiment discontinuously. The distinction matters
when a selected algebra is used to describe a finite observer.

\begin{proposition}[Finite-time operational stability]
\label{prop:operational-stability}
Let $\mathcal L_0,\mathcal L_1$ be finite-dimensional Lindblad generators,
and put $v=\|\mathcal L_1-\mathcal L_0\|_\diamond$. Then
\begin{equation}\label{eq:duhamel-stability}
 \|e^{t\mathcal L_1}-e^{t\mathcal L_0}\|_\diamond
 \le\min\{2,tv\}.
\end{equation}
Consider any experiment with the same initial states, retained ancillas
and permitted instruments in both models. Between its instruments the
system evolves for prescribed intervals $t_1,\ldots,t_m$, with total
$T=\sum_i t_i$. Instruments may depend on earlier outcomes. The total
variation distance between its complete record laws is at most
$\min\{1,Tv/2\}$. Consequently, for a common decision problem with loss
in $[0,L]$, the two optimal risks differ by at most
$L\min\{1,Tv/2\}$. In particular, a Hamiltonian perturbation
$-i\epsilon[K,\cdot]$ gives $v\le2|\epsilon|\|K\|$.
\end{proposition}
\begin{proof}
Duhamel's formula writes the channel difference as
$\int_0^t e^{(t-s)\mathcal L_1}(\mathcal L_1-\mathcal L_0)
e^{s\mathcal L_0}\,ds$. CPTP maps have diamond norm one, giving $tv$;
two channels differ by at most two. A hybrid argument replaces the
evolution intervals one at a time, keeping outcome-controlled
instruments as CPTP maps on the system and classical record. Contractivity
and the triangle inequality give a trace-norm difference at most $Tv$.
The bounded-loss estimate holds uniformly over each common policy, so
it also holds for the infima. Finally,
$\|[K,A]\|_1\le2\|K\|\|A\|_1$ holds with every retained ancilla.
\end{proof}

The estimate needs neither a gap nor a selected algebra. It bounds
distinguishing, storing and acting through the permitted experiment,
not just the state of an isolated system. The available instruments
and the total time remain the same on both sides. A discontinuous
change of permissions, or an unbounded observation time, is a different
comparison.

For an exact noncommuting example, take
\begin{equation}\label{eq:transverse-qubit}
 \mathcal L_\epsilon(\rho)=-i\epsilon[X,\rho]
 +\kappa(Z\rho Z-\rho),\qquad \kappa>0.
\end{equation}
The fixed algebra is $\operatorname{span}\{I,Z\}$ at $\epsilon=0$ and
only $\mathbb CI$ at every $\epsilon\ne0$. To see the latter, the
Hilbert--Schmidt dissipative form of a fixed observable forces
$[Z,O]=0$; the remaining Hamiltonian equation then forces $[X,O]=0$.
Nevertheless, the two initially perfect records $(I\pm Z)/2$ remain
readable for a long time. If $0<2|\epsilon|<\kappa$, put
$\omega=\sqrt{\kappa^2-4\epsilon^2}$. Their Bloch vectors at time $t$
are $\pm(0,y(t),z(t))$, where
\begin{equation}\label{eq:transverse-record}
 z(t)=e^{-\kappa t}\left(\cosh\omega t+
       \frac{\kappa}{\omega}\sinh\omega t\right),\qquad
 y(t)=-\frac{2\epsilon}{\omega}e^{-\kappa t}\sinh\omega t.
\end{equation}
These expressions solve $\dot y=-2\kappa y-2\epsilon z$,
$\dot z=2\epsilon y$, with $(y,z)=(0,1)$ initially. Reading only $Z$
has error $(1-z(t))/2$; allowing the optimal system measurement gives
error $(1-\sqrt{y(t)^2+z(t)^2})/2$. The slow decay rate is exactly
\begin{equation}\label{eq:transverse-rate}
 r=\kappa-\sqrt{\kappa^2-4\epsilon^2}
   =\frac{2\epsilon^2}{\kappa}+O(\epsilon^4/\kappa^3).
\end{equation}
Thus arbitrarily weak transverse coupling erases the infinite-time
record, while its finite-time lifetime diverges as
$\kappa/(2\epsilon^2)$. Exact hulls and fixed algebras describe the
ideal endpoint; finite-time risks describe what remains available on
the way to it. This is also the distinction needed by Foundation's
hulls and the maintenance tasks of T-DOME and Brain Anchor. Neither
the general stability estimate nor this example asserts that every
slow operator space is an algebra or a physical projection.

\section{Records live on the centre}
\label{sec:records}

\subsection{A recording model}

To see what the environment learns we replace the Markovian generator by an exact unitary model. Environmental fragments $E_1,\dots,E_K$ start in states $\omega_k$, and the interaction is
\begin{equation}
U=\sum_{x\in\mathcal X}\proj x\otimes V_1(x)\otimes\cdots\otimes V_K(x),
\label{eq:U}
\end{equation}
with unitaries $V_k(x)$ that depend on the basis state through the coupling, for example $V_k(x)=\exp[-i\theta_k a_k(x)G_k]$ for a fragment generator $G_k$. The joint state after the interaction is
\[
\sigma=\sum_{x,y}\rho_{xy}\ket x\!\bra y\otimes\bigotimes_kV_k(x)\,\omega_k\,V_k(y)^\dagger ,
\]
so the system undergoes the Schur multiplier
\begin{equation}
\rho_{xy}\mapsto c(x,y)\,\rho_{xy},\qquad c(x,y)=\prod_kc_k(x,y),\qquad c_k(x,y)=\Tr\bigl[\omega_kV_k(y)^\dagger V_k(x)\bigr].
\label{eq:coherence}
\end{equation}
The generator \eqref{eq:L} arises from such models in weak-coupling and repeated-collision limits. Here the collision statements are exact for the specified unitary and an initial product $\rho\otimes\bigotimes_k\omega_k$. Fresh independent fragments are an input: recollisions or shared initial fragment correlations do not obey the product formula automatically.

\begin{proposition}[Records carry only the block label]
\label{prop:centre}
Let $\pi$ be a partition of $\mathcal X$ such that every $V_k(x)$ depends on $x$ only through its class $a(x)\in\pi$. Then:
\begin{enumerate}
\item $c(x,y)=1$ whenever $x$ and $y$ lie in one class, so $\E_\pi(\rho)$, and with it the state on every block, is unchanged.
\item For every $O\in\B_\pi$ and every operator $F$ on any collection $\mathcal F$ of fragments,
\[
\Tr[\sigma(O\otimes F)]=\sum_{a\in\pi}\Tr[P_a\rho P_a\,O]\ \Tr[\omega_{\mathcal F}(a)F],\qquad
\omega_{\mathcal F}(a)=\bigotimes_{k\in\mathcal F}V_k(a)\,\omega_k\,V_k(a)^\dagger .
\]
The state of the fragments is $\sum_ap(a)\,\omega_{\mathcal F}(a)$ with $p(a)=\Tr P_a\rho$; it depends on $\rho$ only through the distribution of the label, and given the label, block observables and fragments are uncorrelated.
\end{enumerate}
\end{proposition}

\begin{proof}
(1) $V_k(x)=V_k(y)$ when $x,y$ lie in one class, so $c_k(x,y)=1$. (2) $O_{yx}\ne0$ only when $x$ and $y$ lie in one class $a$, and then the fragment factor is $\omega_{\mathcal F}(a)$ for both. Summing $\rho_{xy}O_{yx}$ over $x,y\in a$ gives $\Tr[P_a\rho P_aO]$.
\end{proof}

Whatever the environment learns under such couplings is the block label, which generates the centre of $\B_\pi$. Many fragments can hold readable copies of that label while leaving the blocks alone. Block dependence alone does not guarantee distinguishability: the conditional states $\omega_k(a)$ must also be different, and a permitted measurement must resolve them. Records and the correlations inside the blocks are thus carried by one selected algebra, the first by its centre and the second by its blocks. That the redundantly available information is classical, and is information about one measurement, holds far more generally \cite{BrandaoPianiHorodecki2015}; a two-qubit system with a decoherence-free subspace that proliferates part of its information while the remaining subspace keeps non-classical features was studied in \cite{Cakmak2021}. What the commuting class adds is the exact location of both kinds of information in one algebra and, next, the price of imperfect symmetry.

\subsection{Leakage costs coherence, and the costs multiply}

\begin{theorem}[Leakage and coherence]
\label{thm:leak}
For basis states $x,y$ and a fragment $k$, let
$D_k(x,y)=\tfrac12\bigl\|V_k(x)\omega_kV_k(x)^\dagger-V_k(y)\omega_kV_k(y)^\dagger\bigr\|_1$
be the fragment's distinguishability of $x$ and $y$, so that an optimal measurement on the fragment identifies one of the two equally likely states with probability $\tfrac12(1+D_k)$ \cite{Helstrom1976}. Then
\begin{equation}
|c_k(x,y)|^2+D_k(x,y)^2\le1,
\label{eq:englert}
\end{equation}
with equality when $\omega_k$ is pure. Consequently
\begin{equation}
|c(x,y)|\le\prod_k\bigl(1-D_k(x,y)^2\bigr)^{1/2}\le\exp\Bigl(-\tfrac12\sum_kD_k(x,y)^2\Bigr),
\label{eq:product}
\end{equation}
and if the coherence between $x$ and $y$ must stay above $1-\varepsilon$, then
\begin{equation}
\sum_kD_k(x,y)^2\le2\ln\frac1{1-\varepsilon}.
\label{eq:budget}
\end{equation}
\end{theorem}

\begin{proof}
Purify $\omega_k=\Tr_R\proj{\Omega}$. Then $c_k(x,y)=\langle\Omega|(V_k(y)^\dagger V_k(x)\otimes I)|\Omega\rangle$ is the overlap of the pure states $(V_k(x)\otimes I)\ket\Omega$ and $(V_k(y)\otimes I)\ket\Omega$, whose trace distance is $(1-|c_k|^2)^{1/2}$. Partial trace over $R$ does not increase trace distance, which gives \eqref{eq:englert}, with equality when there is nothing to trace out. Inequality \eqref{eq:product} follows from \eqref{eq:coherence} and $\ln(1-u)\le-u$, and \eqref{eq:budget} by taking logarithms.
\end{proof}

Inequality \eqref{eq:englert} is Englert's relation between visibility and which-way distinguishability \cite{Englert1996}, applied to one fragment. It bounds what a fragment learns by the coherence it destroys, not the reverse: a fragment in a mixed state can destroy coherence while learning nothing, as in Paper~E's maximally mixed memory example. Purity is one useful preparation resource; mixed fragments can also record when their conditional states differ. For many fragments it becomes the budget \eqref{eq:budget}: the sum of their squared pairwise distinguishabilities is at most $2\ln(1/(1-\varepsilon))$ for a protected pair of basis states. This is not a bound on their total Shannon information about an arbitrary block.

\begin{corollary}[Redundancy against symmetry]
\label{cor:threshold}
Suppose that $R$ fragments each distinguish two states $x,y$ of one block with $D_k(x,y)\ge\beta$. Then $|c(x,y)|\le(1-\beta^2)^{R/2}$, and keeping the coherence above $1-\varepsilon$ requires
\begin{equation}
\beta\le\Bigl(\frac{2\ln(1/(1-\varepsilon))}{R}\Bigr)^{1/2}.
\label{eq:Rhalf}
\end{equation}
In a sequence of models with $R\to\infty$ and $D_k(x,y)=\beta_R=\lambda R^{-\alpha}$, $\alpha>0$, for the $R$ fragments that couple to the block, the coherence tends to zero if $\alpha<\tfrac12$. If the fragment states are pure and no other fragment distinguishes $x$ from $y$, it tends to $1$ for $\alpha>\tfrac12$ and to $e^{-\lambda^2/2}$ for $\alpha=\tfrac12$.
\end{corollary}

\begin{proof}
The bound and \eqref{eq:Rhalf} are \eqref{eq:product} and \eqref{eq:budget}. For pure fragments $|c(x,y)|=(1-\beta_R^2)^{R/2}$ by the equality case, and $\frac R2\ln(1-\lambda^2R^{-2\alpha})\to0,\ -\lambda^2/2,\ -\infty$ for $\alpha>\frac12$, $=\frac12$, $<\frac12$.
\end{proof}

A coupling of the form \eqref{eq:U} is block-constant only by symmetry: the fragments must respond identically to all states of a block. Real couplings break such a symmetry slightly, and each recording fragment then leaks a little. Corollary~\ref{cor:threshold} says that objectivity by redundancy and protection of the blocks can coexist, but only if the symmetry error per fragment falls as $R^{-1/2}$. In the limit of many records this is a sharp threshold in the exponent $\alpha$. It is the closest analogue of a transition in this class. It is a statement about how a symmetry error scales with the number of records, and it is not a phase transition in the thermodynamic sense.

\section{Sector records and protected correlations}
\label{sec:sector}

A mutual-information geometry does not require entanglement. Even local
dephasing preserves
$\rho=\frac12\proj{00}+\frac12\proj{11}$, for which
$I(1:2)=\ln2$. With local $Z$ couplings every selected block is
one-dimensional and every sector state is a product. Its entire edge
comes from sector weights.

For a selected state $\bar\rho=\E_\pi(\rho)=\sum_kp_k\rho_k$, with
$\rho_k=P_k\rho P_k/p_k$, retain the label in a classical register:
\begin{equation}\label{eq:sector-extension}
 \Omega_{SK}=\sum_kp_k\rho_k\otimes\proj{k}_K .
\end{equation}
This extension is physically realized by a perfectly resolving
block-constant recorder followed by measurement or dephasing of its
label register. It is not the coherent source--memory state before
that readout. If the actual fragment states overlap, $K$ is
still a useful mathematical label, but it is not a freely readable
record. Zero-weight sectors are omitted.

\begin{proposition}[The two contributions to correlation]
\label{prop:sector-mi}
For any supplied disjoint subsystems $A,B$ of $S$,
\begin{equation}\label{eq:sector-mi}
 I(A:B)_{\bar\rho}
 =\sum_kp_k I(A:B)_{\rho_k}
  +I(A:K)_\Omega+I(B:K)_\Omega-I(AB:K)_\Omega .
\end{equation}
The last three terms need not have a nonnegative sum.
\end{proposition}
\begin{proof}
For $X=A,B,AB$, the classical--quantum entropy formula gives
$I(X:K)=S(\bar\rho_X)-\sum_kp_kS((\rho_k)_X)$.
Substitution in $S(A)+S(B)-S(AB)$ proves the identity.
\end{proof}

The first term locates correlations protected within sectors; it can
include classical as well as quantum correlations. The correction
describes what mixing sector-conditioned states does to those
correlations. In the locally dephased example above it equals $\ln2$.
In the next recording example it is negative. Thus neither a
noncommutative block nor a classical centre alone determines the
geometry. The sector states, their weights, the site marginals and
the dictionary all enter.

\section{How protected relations are formed}
\label{sec:formation}

A noise commutant specifies what can be stored. It does not determine
whether that content was inherited or created. The distinction already
invalidates a noise-only mereology rule.

\subsection{Disconnected noise can preserve a shared Bell pair}

Take four qubits with only $A_{12}=Z_1+Z_2$ and
$A_{34}=Z_3+Z_4$. Each pair has a zero-magnetization sector with logical
basis $\ket{0_L}=\ket{01}$, $\ket{1_L}=\ket{10}$. The state
\begin{equation}\label{eq:disconnected-bell}
 \ket{\Psi_L}
 =\frac{\ket{0_L}_{12}\ket{0_L}_{34}
             +\ket{1_L}_{12}\ket{1_L}_{34}}{\sqrt2}
\end{equation}
is annihilated by both noise operators and is fixed exactly by
\eqref{eq:L}. Yet $I(12:34)=2\ln2$. The two noise-hypergraph components
are disconnected. Applying a Hadamard to each qubit gives the same
counterexample for the $X$ couplings of Section~\ref{sec:ghz}.
Classical and quantum correlations can both cross disconnected noise
components when they are present initially.

\subsection{A collective recorder forms a flagged protected relation}

Start two qubits in the physical product $\ket{+_X}\ket{+_X}$ and a
fresh three-level memory in $\ket0$. Let $M=Z_1+Z_2$ and let
$P_+,P_0,P_-$ project onto its eigenvalues $2,0,-2$.
For the shift $S_3\ket{k}=\ket{k+1\bmod3}$, use
\begin{equation}\label{eq:formation-recorder}
 U_{\rm rec}=P_+\otimes I+P_0\otimes S_3+P_-\otimes S_3^2 .
\end{equation}
This is a block-constant recording interaction, not a product of
independent local readers. It has a finite Hamiltonian implementation:
choose a Hermitian logarithm $G$ with $e^{-iG}=S_3$ and
$\|G\|=2\pi/3$, put $a(M)=(2I-M)/2$, and apply
$H_{\rm rec}=(\hbar/t_g)a(M)\otimes G$ for a time $t_g>0$.
The same memory couples to both sites. Its preparation, shared physical
port, interaction time, control and subsequent readout are resources.

The complete output is
\begin{equation}\label{eq:formation-pure}
 \ket{\Omega}
 =\tfrac12\ket{00}\ket0
   +\tfrac1{\sqrt2}\ket{\Psi^+}\ket1
   +\tfrac12\ket{11}\ket2,\qquad
 \ket{\Psi^+}=(\ket{01}+\ket{10})/\sqrt2 .
\end{equation}
A computational-basis measurement of the memory produces the
classical--quantum state with probabilities $(1/4,1/2,1/4)$ and
conditional source states
$(\proj{00},\proj{\Psi^+},\proj{11})$. All branches remain in the
protocol. Flag $1$, of probability $1/2$, certifies a Bell pair in the
$M=0$ protected sector; flags $0,2$ certify product states.

If the flag is discarded, the source is instead
\begin{equation}\label{eq:formation-unconditional}
 \rho_*=\tfrac14\proj{00}+\tfrac12\proj{\Psi^+}
                 +\tfrac14\proj{11},\qquad I(1:2)_{\rho_*}=\tfrac12\ln2 .
\end{equation}
It is separable. Indeed it is the uniform collective $Z$-rotation
twirl of $\proj{+_X,+_X}$, a mixture of product rotations.
Its partial-transpose eigenvalues are $0,1/4,1/4,1/2$.
The unconditional source therefore gains a mutual-information
edge without entanglement; separability does not imply that every
correlation measure is classical. The retained flagged ensemble permits a
conditional protected Bell use, with herald probability and every
unsuccessful branch charged. The pure source--memory state
\eqref{eq:formation-pure}, the measured flagged state and the
unconditional source are three different outputs.

For this flagged ensemble,
$\sum_kp_k I(1:2)_{\rho_k}=\ln2$,
$I(1:K)=I(2:K)=\frac12\ln2$ and
$I(12:K)=H(K)=\frac32\ln2$.
Equation~\eqref{eq:sector-mi} recovers the unconditional
$\frac12\ln2$ with correction $-\frac12\ln2$.
Collective dephasing of the product input alone reaches the same
$\rho_*$ but supplies no accessible flag. A flag requires the
specified recording interaction and actual readout.

\subsection{A real link forms a deterministic protected Bell pair}

The distinction is sharper in the four-qubit example. Define
\[
 L^x_{ab}=\tfrac12(X_aX_b+Y_aY_b),\qquad
 H_{\rm link}=\hbar J L^x_{12}L^x_{34},\quad J>0 .
\]
$L^x_{ab}$ exchanges $\ket{01}$ and $\ket{10}$ and vanishes on
$\ket{00},\ket{11}$. Consequently
$[H_{\rm link},A_{12}]=[H_{\rm link},A_{34}]=0$.
Starting from the physical product $\ket{01}_{12}\ket{01}_{34}$,
the full noisy evolution remains in the joint zero sector and is
exactly
\begin{equation}\label{eq:formed-logical-bell}
 \ket{\psi(t)}
 =\cos(Jt)\ket{0_L0_L}-i\sin(Jt)\ket{1_L1_L}.
\end{equation}
At $t=\pi/(4J)$ it has $I(12:34)=2\ln2$ and one ebit across
the pair partition. The noise term vanishes throughout the trajectory.

This is new protected correlation from an initially uncorrelated
state. It uses an actual cross-pair Hamiltonian link, here a four-body
control term. The noise graph remains disconnected but the complete
interaction graph does not. The factorization, prepared input, gate
permissions, coupling strength and duration are supplied; there is
no inferred free implementation or heat price for the link.

\subsection{The process constraint and learned use}

\begin{proposition}[Correlation provenance across independent components]
\label{prop:provenance}
Let a supplied partition of a carrier have initially independent
states and auxiliary resources. Suppose its full channel, including
record registers, factors across these components. Then it creates
no cross-component mutual information. If the input is correlated,
a product channel cannot increase the mutual information across any
bipartition of the components, but it can preserve initially shared
classical or quantum correlations.
\end{proposition}
\begin{proof}
A product channel sends a product input to a product output, including
its local records. For a correlated input,
$I(A:B)=D(\rho_{AB}\|\rho_A\otimes\rho_B)$, so the second statement
is data processing under the product channel.
\end{proof}

An implementation satisfies the factor premise when its Hamiltonian
and noise generators act separately on the components, its fresh
ports are independent, and every control and instrument is
component-local. Shared random controls, communication or common
memories must be added to the process description: they can correlate
unconditional outputs. Conditioning on a joint event can also change
correlations and is not a product unconditional channel. The
proposition is a diagnostic, not a claim that every process with a
disconnected noise graph factors.

The constructions above do more than add initial-state qualifications
to a conjecture. They exhibit how the missing relations can form:
a collective port records an indistinguishability sector, while an
interaction inside protected sectors builds usable quantum content.
The other indispensable step is learning which encoding an unknown
physical plant protects. Foundation~III \cite{PaperFIII} implements
that step for $A_s=n_1+s n_2$ with unknown $s=\pm1$. A diagnostic
sharing the plant's actual sign writes a record; that record, not
$s$, selects odd or even encoding and the matching decoder. The same
logical content and its old reference survive from before diagnosis.
The record-controlled frame changes the later contact, rather than
merely changing a description of the fixed commutant.

That result assumes a physical relation between diagnostic and storage
couplings. A diagnostic carrying an unrelated sign would not learn
the protected frame. Its retained-history scalar/JT implementation
also contains real body forces, a source and a second acquisition;
it is not identified with the full quantum material model.
The complementary contribution here is the separation of formation,
protection and use: each can fail while the others hold. A noiseless
block can be unreadable, a readable label can outlive its quantum
content, and a learned encoder can protect content without supplying
every logical gate.

\section{A geometry selected by the coupling}
\label{sec:ghz}

We now apply these results to the state and the reconstruction rule of Paper~A. The rule acts on the complete graph of the sites: the edge $\{i,j\}$ has length $1/I_{ij}$, where $I_{ij}$ is the mutual information of sites $i$ and $j$, and distances are shortest-path lengths. We extend it to $I_{ij}=0$ by giving that edge infinite length, so that a site whose edges all have zero mutual information is disconnected. The state is the three-qubit GHZ state $(\ket{000}+\ket{111})/\sqrt2$. In the eigenbasis $\ket\pm$ of the Pauli operator $X$ it reads
\[
\ket{\mathrm{GHZ}}=\tfrac12\bigl(\ket{+++}+\ket{+--}+\ket{-+-}+\ket{--+}\bigr).
\]
All couplings below are sums of $X_i$, hence diagonal in this basis, and Sections~\ref{sec:selected}--\ref{sec:records} apply. A local coupling $A=X_i$ with rate $\gamma_i$ is the bit-flip noise of Paper~A, with flip probability $p_i=(1-e^{-2\gamma_it})/2$. A collective coupling $A=X_1+X_2+X_3$ dephases between different values of the $X$-magnetization $M=\sum_ix_i$ and leaves states of equal $M$ coherent.

\begin{proposition}[Three couplings, three geometries]
\label{prop:three}
Let $\Psi^+=(\ket{+-}+\ket{-+})/\sqrt2$ and $\tilde W=(\ket{+--}+\ket{-+-}+\ket{--+})/\sqrt3$. The long-time selected states $\E_{\pi_\infty}(\proj{\mathrm{GHZ}})$ and their geometries are:
\begin{enumerate}
\item Collective coupling $X_1+X_2+X_3$: $\tfrac14\proj{+++}+\tfrac34\proj{\tilde W}$. Each pair has $I_{ij}=\tfrac12\ln2$; the geometry is an equilateral triangle with side $2/\ln2=2.8854$.
\item Collective coupling $X_1+X_2$ and local coupling $X_3$: $\tfrac14\proj{+++}+\tfrac12\proj{\Psi^+}_{12}\otimes\proj{-}_3+\tfrac14\proj{--+}$. Here $I_{12}=\tfrac12\ln2$ and $I_{13}=I_{23}=0$; the geometry is a single edge, and site $3$ is disconnected.
\item Local couplings $X_1,X_2,X_3$: the uniform mixture of the four even strings $\ket{+++},\ket{+--},\ket{-+-},\ket{--+}$. Every $I_{ij}=0$; no edge survives.
\end{enumerate}
\end{proposition}

\begin{proof}
The pinching keeps the amplitudes inside each class of equal eigenvalue tuples. (1) The classes met by the GHZ state are $M=3$ and $M=-1$. The marginal of sites $1,2$ is $\frac14\proj{++}+\frac12\proj{\Psi^+}+\frac14\proj{--}$, with eigenvalues $\frac14,\frac12,\frac14$ and entropy $\frac32\ln2$; each single-site marginal is $I/2$; so $I_{12}=2\ln2-\frac32\ln2$, and the other pairs follow by symmetry. (2) The classes are labelled by $(x_1+x_2,x_3)$; the strings $\ket{+--}$ and $\ket{-+-}$ share the label $(0,-)$. The marginal of sites $1,2$ is as in (1). The marginal of sites $1,3$ is uniform on the four strings $\ket{++},\ket{+-},\ket{--},\ket{-+}$, and so is that of sites $2,3$. (3) Every pair marginal of the uniform mixture of even strings is uniform. Direct evaluation of the full eight-dimensional density matrices gives the same values.
\end{proof}

The state and the rule are the same in all three cases; only the coupling differs. At zero time each pair has $I_{ij}=\ln2$, Paper~A's GHZ geometry at zero noise. Under local noise Paper~A's triangles relax toward the empty graph; a collective coupling stops the relaxation at an equilateral triangle, and a coupling that is collective on one pair keeps exactly that edge. The surviving correlations come from blocks of dimension greater than one: the three-dimensional block $M=-1$, which contains the entangled state $\tilde W$, in case (1), and the two-dimensional block spanned by $\ket{+--}$ and $\ket{-+-}$ in case (2). In case (3) the selected algebra is commutative with product minimal projections, every selected state is classically correlated, and for this state the classical pair correlations vanish as well. By Proposition~\ref{prop:selected}(4) the channels at any finite time generate the full matrix algebra in all three cases, so an algebra label cannot tell them apart.

\subsection{The lifetime of the selected geometry}

Real couplings are not exactly collective. Add weak local couplings $X_i$ with rates $\gamma_i$ to the collective coupling with rate $\gamma_c$. The exact fixed algebra is then commutative and the long-time geometry is empty, but on a window the magnetization blocks are still selected. With $\pi$ the partition by $M$, the rates \eqref{eq:rates} are $\Gin(\pi)=2(\gamma_1+\gamma_2+\gamma_3)$ and $\Gout(\pi)=2\gamma_c+2\min_i\gamma_i$, while the coherences of the GHZ state between its two blocks decay at the rates $8\gamma_c+2(\gamma_i+\gamma_j)$.

\begin{lemma}[Plateau bound]
\label{lem:plateau}
For a state $\rho$ and a partition $\pi$, let $\Gout(\rho,\pi)$ be the smallest $D(x,y)$ over pairs in different classes with $\rho_{xy}\ne0$, and let $r$ be the number of basis vectors with $\rho_{xx}\ne0$. Then for all $t\ge0$
\begin{equation}
\|\Phi_t(\rho)-\E_\pi(\rho)\|_1\le\sqrt r\,e^{-t\Gout(\rho,\pi)}\,\|\rho-\E_\pi(\rho)\|_2+t\,\Gin(\pi).
\label{eq:plateau}
\end{equation}
\end{lemma}

\begin{proof}
Split $\Phi_t(\rho)-\E_\pi\rho=\Phi_t(\rho-\E_\pi\rho)+(\Phi_t(\E_\pi\rho)-\E_\pi\rho)$. The second term is bounded by Theorem~\ref{thm:window}(1). The first is the Schur product of $\rho-\E_\pi\rho$ with multipliers of modulus at most $e^{-t\Gout(\rho,\pi)}$; its rank is at most $r$, because a row of a positive matrix with zero diagonal entry vanishes, and $\|Z\|_1\le\sqrt{\operatorname{rank}Z}\,\|Z\|_2$.
\end{proof}

The bound turns into a certified geometry through continuity of entropy. If two states of a qubit pair are at trace distance $\delta=\tfrac12\|\rho-\sigma\|_1\le\frac12$, then by the Fannes--Audenaert inequality \cite{Audenaert2007} and Winter's bound for conditional entropy \cite{Winter2016} their mutual informations differ by at most
\begin{equation}
f(\delta)=h(\delta)+2\delta\ln2+(1+\delta)\,h\Bigl(\frac{\delta}{1+\delta}\Bigr),
\label{eq:fdelta}
\end{equation}
with $h$ the binary entropy in nats. Partial traces do not increase trace distance, and Paper~A's stability bound for graph metrics converts a uniform change $f<m$ of mutual informations bounded below by $m$ into a change of at most $2f/(m(m-f))$ of every distance among three sites.

\begin{table}[ht]
\centering
\begin{tabular}{rccc}
\toprule
$t$ & $\|\Phi_t(\rho)-\E_\pi\rho\|_1$ & bound \eqref{eq:plateau} & distances $(d_{12},d_{13},d_{23})$\\
\midrule
$0$ & $0.8660$ & $1.2247$ & $(1.4427,\ 1.4427,\ 1.4427)$\\
$1$ & $0.0010$ & $0.0018$ & $(2.8968,\ 2.9034,\ 2.9066)$\\
$10$ & $0.0094$ & $0.0140$ & $(2.9728,\ 3.0221,\ 3.0458)$\\
$100$ & $0.0897$ & $0.1400$ & $(3.5620,\ 3.9855,\ 4.2034)$\\
$300$ & $0.2442$ & $0.4200$ & $(4.8943,\ 6.5032,\ 7.4508)$\\
$1000$ & $0.5969$ & $1.4000$ & $(12.559,\ 28.867,\ 41.426)$\\
\bottomrule
\end{tabular}
\caption{GHZ state under a collective coupling with $\gamma_c=1$ and local couplings with $(\gamma_1,\gamma_2,\gamma_3)=10^{-4}(1,2,4)$. The plateau value of every distance is $2/\ln2=2.8854$.}
\label{tab:plateau}
\end{table}

\begin{proposition}[A certified plateau]
\label{prop:plateau}
For $\gamma_c=1$ and $(\gamma_1,\gamma_2,\gamma_3)=10^{-4}(1,2,4)$, the bounds \eqref{eq:plateau} and \eqref{eq:fdelta} certify that every distance of the GHZ geometry lies within $25\%$ of the plateau value $2/\ln2$ for $0.716\le t\le3.58$. Numerical root finding of the exact evolution gives the approximate $25\%$ window $0.053\lesssim t\lesssim53.6$ and $5\%$ window $0.163\lesssim t\lesssim8.86$. If the local rates are scaled to $s(1,2,4)$ with $s\le10^{-3}$, the late edges scale to leading order as $1/s$: approximately $5.36\times10^{-3}/s$ for the $25\%$ window, $8.86\times10^{-4}/s$ for the $5\%$ window and, when $s\le10^{-4}$ so that it exists, $3.58\times10^{-4}/s$ for the certified $25\%$ window. The windows begin at times of order $1/\gamma_c$.
\end{proposition}

\begin{proof}
For this state $r=4$, $\|\rho-\E_\pi\rho\|_2=\sqrt6/4$, $\Gout(\rho,\pi)=8.0006$ and $\Gin(\pi)=1.4\times10^{-3}$. The right side of \eqref{eq:plateau} is the convex function $b(t)=(\sqrt6/2)e^{-8.0006t}+0.0014t$. On $[0.716,3.58]$ its maximum is at an endpoint and is less than $0.0050121$. The trace distance used in~\eqref{eq:fdelta} is
$\delta(t)=\tfrac12\|\Phi_t(\rho)-\E_\pi\rho\|_1\le b(t)/2$,
not $b(t)$. The increasing entropy and graph bounds therefore give a distance error less than $0.72130<\frac14\cdot2/\ln2$ throughout that interval. The exact-evolution window edges were located by root finding of the Schur multipliers. Once the coherences between blocks have decayed, the in-block multipliers depend on $t$ only through $st$, which gives the late-edge scaling. The root values are numerical diagnostics, distinct from the conservative analytic interval.
\end{proof}

Table~\ref{tab:plateau} shows the three stages. The GHZ geometry, with side $1/\ln2$, relaxes on the time scale $1/\gamma_c$ to the selected triangle with side $2/\ln2$; the triangle persists while $t\Gin(\pi)$ is small; then the unequal local rates deform it into scalene triangles, which are not isometric to the plateau, before every edge diverges. This is Paper~A's non-isometry produced by one time-ordered dynamics rather than by a comparison of separately chosen channels. Proposition~\ref{prop:plateau} gives a state-specific stability certificate for Paper~A's declared graph dictionary on an explicit window. The plateau exists when its end lies after its beginning, that is, when the ratio of in-block to between-block rates is small, as Theorem~\ref{thm:window} requires.

The constants in Proposition~\ref{prop:plateau} are small. At the end of the exact $25\%$ window the three in-block coherences have fallen by only $3.2$, $5.2$ and $6.2$ per cent. The selected pair marginal has a zero eigenvalue, and in-block decoherence populates it, so pair entropies respond to a coherence loss $\delta$ like $\delta\ln(1/\delta)$. Geometries read from entropies are therefore far more fragile than the coherences that carry them, and the gap between the certified and the exact window comes mostly from this logarithmic response together with worst-case continuity bounds.

\subsection{Imperfect redundant records detach a vertex}

Finally we let fragments record the magnetization with a small asymmetry. Each of $K$ fragment qubits starts in $\ket0$ and interacts through $V_k(x)=\exp[-i\theta(M(x)+\eta x_1)Y]$ with $\theta=\pi/8$, where $Y$ is the Pauli operator of the fragment. For $\eta=0$ the coupling is block-constant. The two magnetizations present in the GHZ state, $M=3$ and $M=-1$, then produce orthogonal fragment states, so every fragment holds a perfect copy of the label, and by Proposition~\ref{prop:centre} the equilateral triangle survives for any number of fragments. For $\eta>0$ each fragment distinguishes strings of the block $M=-1$ that differ at site $1$, with $D_k=|\sin2\theta\eta|$.

\begin{table}[ht]
\centering
\begin{tabular}{cccccc}
\toprule
$\eta$ & $K$ & coherence of $\ket{+--},\ket{-+-}$ & $I_{12}=I_{13}$ & $I_{23}$ & label error\\
\midrule
$0.02$ & $1$ & $0.99988$ & $0.34631$ & $0.34657$ & $6.2\times10^{-5}$\\
$0.02$ & $100$ & $0.98774$ & $0.32790$ & $0.34657$ & $6.2\times10^{-5}$\\
$0.02$ & $1000$ & $0.88393$ & $0.23581$ & $0.34657$ & $6.2\times10^{-5}$\\
$0.05$ & $1$ & $0.99923$ & $0.34525$ & $0.34657$ & $3.9\times10^{-4}$\\
$0.05$ & $100$ & $0.92577$ & $0.26724$ & $0.34657$ & $3.9\times10^{-4}$\\
$0.05$ & $1000$ & $0.46243$ & $0.05555$ & $0.34657$ & $3.9\times10^{-4}$\\
\bottomrule
\end{tabular}
\caption{Recording the magnetization with a site-$1$ asymmetry $\eta$. The last column is the minimal error probability with which a measurement on one fragment identifies the label $M$, whose prior is $(\frac14,\frac34)$. Coherences are relative to their initial values; the coherence between $\ket{-+-}$ and $\ket{--+}$ stays equal to $1$.}
\label{tab:records}
\end{table}

Table~\ref{tab:records} shows the effect. Every fragment is a nearly perfect and readable record of the label for both asymmetries, so the loss of geometry is not a price paid for better records. It is the small leakage about site $1$, multiplied over fragments as in Theorem~\ref{thm:leak}. The coherences between strings that differ at site $1$ decay as $(1-D_k^2)^{K/2}$, the edges to site $1$ weaken, and as $K\to\infty$ the state approaches $\frac14\proj{+++}+\frac14\proj{+--}+\frac12\proj{-}_1\otimes\proj{\Psi^+}_{23}$. Then $I_{12}=I_{13}=0$ and $I_{23}=\frac12\ln2$: site $1$ is detached and the geometry becomes the single edge of Proposition~\ref{prop:three}(2), now on sites $2,3$. With $\eta_K=0.5K^{-\alpha}$ the coherence after $K=10^6$ fragments is $0$ for $\alpha=\frac14$, $0.9258=e^{-(\pi/8)^2/2}$ for $\alpha=\frac12$, and $0.9999$ for $\alpha=\frac34$, as Corollary~\ref{cor:threshold} predicts with $\lambda=2\theta\cdot0.5$ to leading order.

\section{What selection settles}
\label{sec:meaning}

The selected partition is fixed by the supplied dynamics, window and
tolerance, and is robust to perturbations preserving the rate gap.
This differs in emphasis from QBism's state assignments and relational
quantum mechanics' reference-system dependence
\cite{Fuchs2014,Rovelli1996}, without deciding their wider interpretations.
Public agreement further needs actual fragments and calibrated readers;
Paper~A provides one finite realization. A stable algebra alone does not
grant those permissions.

The GHZ comparison changes a state under one graph dictionary, not the
coordinates of an unchanged metric. Sector mixing can sustain edges
after quantum content is lost; the leaking GHZ block instead supports
the specific edges destroyed in Table~\ref{tab:records}. Formation
requires the collective port or Hamiltonian link of
Section~\ref{sec:formation}, not noise selection alone.
Removing either resource supplies its corresponding opposing control.

AQFT and Consciousness~I start with their local or current-net carriers
\cite{Haag1996,PaperCI}. Observable-induced factorization and quantum
mereology address when retained relations support an effective subsystem
description \cite{Zanardi2004,Cotler2019,CarrollSingh2021}. Decoherence and
Everettian accounts share the underlying dynamics
\cite{Zurek2003,Wallace2012}; the issue here is whether its resulting
structure is readable, operable and persistent for a finite task.

\section{Beyond the commuting class}
\label{sec:conj}

\begin{conjecture}[Process-selected effective subsystems]
\label{conj:mereology}
On an initially specified carrier, robust effective subsystems form
through a history of interactions that both builds relations and
makes a noncommuting set of internal operations reproducible across
a finite task window. The relevant connectivity is that of the
complete physical process, including recording ports and
Hamiltonian controls, not the collective-noise hypergraph alone.
Within a separated-rate regime, a learned encoding can turn the
process's indistinguishable sectors into stable operational
components; their boundaries change when the reproducible
operations or their lifetime change.
\end{conjecture}

The conjecture concerns effective factorization, not the existence
of elementary sites. It predicts more than correlation alone:
successful formation should yield an encoder, decoder and
repeatable operations on the formed component, with calibration
records and finite error bounds. The exact formation examples give
protected relational content; Foundation~III adds record-controlled
encoding of unknown content. Together they support a route from
coupling through history to changed capability. They do not yet
supply a general autonomous acquisition of a complete internal gate
set. A proposed component that has only a state-dependent correlation
but no reproducible internal operations would fail this criterion.
The disconnected Bell counterexample rules out noise-only
connectivity as a necessary condition, while
Proposition~\ref{prop:provenance} identifies the real no-creation
obstruction for product processes.

For general Lindblad dynamics, a separated spectrum defines slow
states and an effective metastable evolution
\cite{Macieszczak2016}. A finite-window extension should select the
operational structure of this manifold rather than assume that slow
observables close under multiplication. Even complete positivity
and idempotence do not ensure closure in the original multiplication.
For example,
\begin{equation}\label{eq:cp-range}
 P(X)=\operatorname{diag}\bigl(x_{11},x_{22},
                                      (x_{11}+x_{22})/2\bigr)
\end{equation}
is a unital completely positive idempotent on $M_3$.
It is completely positive because it is a sum of two positive
evaluation maps into positive diagonal matrices. Its range contains
$E=\operatorname{diag}(1,0,1/2)$ but not
$E^2=\operatorname{diag}(1,0,1/4)$.
A completely positive projection gives its range a Choi--Effros
product $x\circ y=P(xy)$ \cite{ChoiEffros1977,BlecherNeal2016}.
This abstract product does not make $xy$ a physically permitted
observable or control in the original representation. An ordinary
conditional expectation onto a represented subalgebra needs the
additional algebraic closure. A metastable spectral projection may
not even be completely positive. Thus a general selection theory
must provide actual instruments or recovery maps, their errors and
their operational closure, rather than infer them from a slow
spectral space.

The physical price of that capability is the next link to Q-RAIF
and the gravitational line. A passive GHZ geometry has lifetime
of order $1/s$ under symmetry breaking, while independent
redundant fragments require an $R^{-1/2}$ protecting symmetry.
Q-RAIF~B realizes correction as field coupling with work and stress
\cite{Liu2026QRAIFB}; HAFF~D distinguishes entropy removal,
usable side information and supplied purity, and gives specified
source--response constructions \cite{Liu2026PaperD}.
These results suggest a maintained-structure mechanism:
actual records learn a protected frame, operations install it,
maintenance changes the physical source, and the response changes
subsequent records. Foundation~III realizes a finite retained-history
scalar/JT return, whereas the full quantum material source requires
its own support and preparation account. Connecting those
implementations and renewing their resources is a substantive
requirement, not an irreducible heat flux inferred from algebra
inclusion.

Many-site and continuum selection require further work when rate
gaps close. Consciousness~I's ideal fixed-net closure and finite-time
noise estimates are not already a continuum uniqueness theorem for
window-selected blocks. Likewise, the graph dictionary is supplied
here. A general emergence of spacetime needs a physical family of
readers and a reconstruction law that remains controlled as the
operational structure grows.

\section{Conclusion}
\label{sec:conclusion}

A commuting coupling determines its fixed block algebra and its
long-time pinching. On a separated finite window it selects a
unique block partition in the supplied joint eigenbasis; the
selected algebra contains every stable block algebra of that form.
The actual finite-time channel has full operator range.
Block-constant record couplings report sector weights without
resolving protected content. Independent leakage multiplies the
coherence loss and imposes a sharp redundancy--symmetry threshold.

The relation to geometry has two parts. Sector mixing and
within-sector content contribute through the exact
identity~\eqref{eq:sector-mi}; neither alone determines the mutual
information graph. For the GHZ source, collective, pair-collective
and local couplings yield a triangle, one edge and no edge under
one rule. Weak asymmetry gives a finite plateau and then destroys
its supporting edges. Noncommuting perturbations can erase the
exact fixed algebra without abruptly erasing a finite record or
task capability.

Formation requires a history. An initially shared logical Bell
pair can survive between noise-disconnected protected components,
so a noise graph cannot delimit all surviving quantum relations.
A collective recorder forms a flagged protected Bell branch from
a product input, with all branches retained; an actual Hamiltonian
link forms a deterministic Bell pair protected throughout its
generation. These constructions reduce the missing formation step
to identifiable physical operations without hiding their costs.
Record-controlled encoding supplies the further route from learned
relations to changed capability. The remaining general problem is
to derive reproducible subsystem operations and their sustained
maintenance from such histories, rather than merely rename
correlated states as subsystems.

\begin{thebibliography}{99}

\bibitem{Liu2026PaperA}
S.~Liu, \emph{Physical Access and Correlation Geometry: The Holographic Alaya--Field Framework}, 2026; series record DOI: 10.5281/zenodo.18361706.

\bibitem{Liu2026PaperD}
S.~Liu, \emph{Selection, Maintenance and Gravitational Response: Modular Constraints and a Physical Return}, (2026), series record DOI: 10.5281/zenodo.18388881.

\bibitem{Liu2026PaperE}
S.~Liu, \emph{Measurement as Accessibility: Measurement Permissions, Stable Records and Public Readout}, (2026), series record DOI: 10.5281/zenodo.18400065.

\bibitem{PaperCI}
S.~Liu, \emph{Geometry Relative to a Cut: Selection, Generation and Modular Inseparability}, (2026), series record DOI: 10.5281/zenodo.19607387.

\bibitem{Liu2026QRAIFB}
S.~Liu, \emph{Thermodynamic Stability Constraints on Persistent Open Quantum Subsystems: Control Algebra, Physical Actuation and Energy Cost}, (2026), series record DOI: 10.5281/zenodo.18525890.

\bibitem{PaperFIII}
S.~Liu, \emph{Foundation III: The Perfuming Loop. Held Content, Learned Frames and Physical Return}, 2026, in the Genesis Trilogy working corpus.

\bibitem{ChoiEffros1977}
M.-D.~Choi and E.~G.~Effros, \emph{Injectivity and operator spaces},
J.~Funct.~Anal.\ \textbf{24}, 156--209 (1977).

\bibitem{BlecherNeal2016}
D.~P.~Blecher and M.~Neal, \emph{Completely contractive projections on operator algebras},
Pacific J.~Math.\ \textbf{283}, 289--324 (2016).

\bibitem{Zurek1982}
W.~H.~Zurek, \emph{Environment-induced superselection rules}, Phys.\ Rev.\ D \textbf{26}, 1862 (1982).

\bibitem{Zurek2003}
W.~H.~Zurek, \emph{Decoherence, einselection, and the quantum origins of the classical}, Rev.\ Mod.\ Phys.\ \textbf{75}, 715 (2003).

\bibitem{Zurek2009}
W.~H.~Zurek, \emph{Quantum Darwinism}, Nature Physics \textbf{5}, 181 (2009).

\bibitem{ZanardiRasetti1997}
P.~Zanardi and M.~Rasetti, \emph{Noiseless quantum codes}, Phys.\ Rev.\ Lett.\ \textbf{79}, 3306 (1997).

\bibitem{KnillLaflammeViola2000}
E.~Knill, R.~Laflamme, and L.~Viola, \emph{Theory of quantum error correction for general noise}, Phys.\ Rev.\ Lett.\ \textbf{84}, 2525 (2000).

\bibitem{Frigerio1978}
A.~Frigerio, \emph{Stationary states of quantum dynamical semigroups}, Commun.\ Math.\ Phys.\ \textbf{63}, 269 (1978).

\bibitem{BlumeKohout2010}
R.~Blume-Kohout, H.~K.~Ng, D.~Poulin, and L.~Viola, \emph{Information-preserving structures: A general framework for quantum zero-error information}, Phys.\ Rev.\ A \textbf{82}, 062306 (2010).

\bibitem{FacchiPascazio2002}
P.~Facchi and S.~Pascazio, \emph{Quantum Zeno subspaces}, Phys.\ Rev.\ Lett.\ \textbf{89}, 080401 (2002).

\bibitem{Macieszczak2016}
K.~Macieszczak, M.~Gu\c{t}\u{a}, I.~Lesanovsky, and J.~P.~Garrahan, \emph{Towards a theory of metastability in open quantum dynamics}, Phys.\ Rev.\ Lett.\ \textbf{116}, 240404 (2016).

\bibitem{BrandaoPianiHorodecki2015}
F.~G.~S.~L.~Brand\~ao, M.~Piani, and P.~Horodecki, \emph{Generic emergence of classical features in quantum Darwinism}, Nature Commun.\ \textbf{6}, 7908 (2015).

\bibitem{Cakmak2021}
B.~\c{C}akmak, \"O.~E.~M\"ustecapl\i{}o\u{g}lu, M.~Paternostro, B.~Vacchini, and S.~Campbell, \emph{Quantum Darwinism in a composite system: Objectivity versus classicality}, Entropy \textbf{23}, 995 (2021).

\bibitem{Helstrom1976}
C.~W.~Helstrom, \emph{Quantum Detection and Estimation Theory}, Academic Press (1976).

\bibitem{Englert1996}
B.-G.~Englert, \emph{Fringe visibility and which-way information: An inequality}, Phys.\ Rev.\ Lett.\ \textbf{77}, 2154 (1996).

\bibitem{Audenaert2007}
K.~M.~R.~Audenaert, \emph{A sharp continuity estimate for the von Neumann entropy}, J.\ Phys.\ A \textbf{40}, 8127 (2007).

\bibitem{Winter2016}
A.~Winter, \emph{Tight uniform continuity bounds for quantum entropies: Conditional entropy, relative entropy distance and energy constraints}, Commun.\ Math.\ Phys.\ \textbf{347}, 291 (2016).

\bibitem{Haag1996}
R.~Haag, \emph{Local Quantum Physics: Fields, Particles, Algebras}, Springer-Verlag (1996).

\bibitem{Rovelli1996}
C.~Rovelli, \emph{Relational quantum mechanics}, Int.\ J.\ Theor.\ Phys.\ \textbf{35}, 1637 (1996).

\bibitem{Fuchs2014}
C.~A.~Fuchs, N.~D.~Mermin, and R.~Schack, \emph{An introduction to QBism with an application to the locality of quantum mechanics}, Am.\ J.\ Phys.\ \textbf{82}, 749 (2014).

\bibitem{Wallace2012}
D.~Wallace, \emph{The Emergent Multiverse}, Oxford University Press (2012).

\bibitem{Zanardi2004}
P.~Zanardi, D.~A.~Lidar, and S.~Lloyd, \emph{Quantum tensor product structures are observable induced}, Phys.\ Rev.\ Lett.\ \textbf{92}, 060402 (2004).

\bibitem{Cotler2019}
J.~Cotler, G.~Penington, and D.~H.~Ranard, \emph{Locality from the spectrum}, Commun.\ Math.\ Phys.\ \textbf{368}, 1267 (2019).

\bibitem{CarrollSingh2021}
S.~M.~Carroll and A.~Singh, \emph{Quantum mereology: Factorizing Hilbert space into subsystems with quasiclassical dynamics}, Phys.\ Rev.\ A \textbf{103}, 022213 (2021).

\end{thebibliography}

\end{document}

